%% Đề 39 -- de-tuyen-sinh-lop-10-mon-toan-nam-2026-2027-so-gddt-lang-son.pdf
\deheader{ĐỀ SỐ 39}%
  {Kỳ thi Tuyển sinh vào lớp 10 THPT — Năm học 2026--2027}%
  {Sở GD\&ĐT Lạng Sơn | Chỉ phần Đại số}
\addcontentsline{toc}{section}{Đề 39 — Tuyển sinh vào lớp 10 THPT — Năm học 2026--2027}

\begin{de_cau}{Câu 1}{2,0 điểm}
\begin{enumerate}[label=\alph*)]
\item Tính giá trị của biểu thức: $A = 3\sqrt{2} - \sqrt{18} + \sqrt{25}$.
\item Rút gọn biểu thức: $B = \left(\dfrac{\sqrt{x}}{\sqrt{x}-1} - \dfrac{1}{x-\sqrt{x}}\right) : \dfrac{\sqrt{x}+1}{\sqrt{x}-1}$ với $x > 0,\; x \neq 1$.
\end{enumerate}
\end{de_cau}

\begin{loi_giai_cau}
\textbf{a)} Ta có $\sqrt{18} = 3\sqrt{2}$ và $\sqrt{25} = 5$. Do đó:
\[A = 3\sqrt{2} - 3\sqrt{2} + 5 = \boxed{5}.\]

\textbf{b)} Với $x > 0, x \neq 1$:
\[B = \left(\frac{x}{\sqrt{x}(\sqrt{x}-1)} - \frac{1}{\sqrt{x}(\sqrt{x}-1)}\right) \cdot \frac{\sqrt{x}-1}{\sqrt{x}+1} = \frac{x-1}{\sqrt{x}(\sqrt{x}-1)} \cdot \frac{\sqrt{x}-1}{\sqrt{x}+1} = \frac{(\sqrt{x}-1)(\sqrt{x}+1)}{\sqrt{x}(\sqrt{x}+1)} = \boxed{\frac{\sqrt{x}-1}{\sqrt{x}}}.\]
\end{loi_giai_cau}

\begin{de_cau}{Câu 2}{2,0 điểm}
\begin{enumerate}[label=\alph*)]
\item Giải hệ phương trình: 
\[ \begin{cases} 2x + y = 5 \\ 3x - 2y = 4 \end{cases} \]
\item Cho phương trình: $x^2 - 2(m-1)x + 2m - 5 = 0$ ($m$ là tham số). Tìm $m$ để phương trình có hai nghiệm phân biệt $x_1, x_2$ thỏa mãn: $x_1^2 + x_2^2 = 10$.
\end{enumerate}
\end{de_cau}

\begin{loi_giai_cau}
\textbf{a)} Nhân phương trình đầu với 2 rồi cộng với phương trình thứ hai:
\[4x + 2y + 3x - 2y = 10 + 4 \iff 7x = 14 \iff x = 2.\]
Thay $x = 2$ vào phương trình đầu: $2(2) + y = 5 \implies y = 1$.
Vậy nghiệm của hệ là $\boxed{(x; y) = (2; 1)}$.

\textbf{b)} Ta có:
\[\Delta' = [-(m-1)]^2 - (2m-5) = m^2 - 2m + 1 - 2m + 5 = m^2 - 4m + 6 = (m-2)^2 + 2 > 0 \quad \forall m.\]
Phương trình luôn có hai nghiệm phân biệt với mọi $m$.
Theo Vi-ét: $x_1+x_2 = 2(m-1)$ và $x_1 x_2 = 2m-5$.
Ta có:
\[x_1^2 + x_2^2 = (x_1+x_2)^2 - 2x_1 x_2 = 4(m-1)^2 - 2(2m-5) = 4m^2 - 12m + 14.\]
Để $x_1^2 + x_2^2 = 10$:
\[4m^2 - 12m + 14 = 10 \iff 4m^2 - 12m + 4 = 0 \iff m^2 - 3m + 1 = 0.\]
Giải phương trình ta được:
\[\boxed{m = \dfrac{3 \pm \sqrt{5}}{2}}.\]
\end{loi_giai_cau}

\begin{de_cau}{Câu 3}{1,5 điểm}
Một mảnh vườn hình chữ nhật có diện tích $360\text{ m}^2$. Nếu tăng chiều rộng thêm $2\text{ m}$ và giảm chiều dài đi $6\text{ m}$ thì diện tích mảnh vườn không thay đổi. Tính chu vi ban đầu của mảnh vườn.
\end{de_cau}

\begin{loi_giai_cau}
Gọi chiều rộng ban đầu là $x\text{ (m)}$ ($x > 0$).
Chiều dài ban đầu là $\dfrac{360}{x}\text{ (m)}$ (điều kiện $\dfrac{360}{x} > 6 \implies x < 60$).
Khi tăng chiều rộng thêm $2\text{ m}$ và giảm chiều dài $6\text{ m}$, ta có phương trình:
\[(x + 2)\left(\frac{360}{x} - 6\right) = 360 \iff 360 - 6x + \frac{720}{x} - 12 = 360 \iff -6x - 12 + \frac{720}{x} = 0.\]
Nhân cả hai vế với $-\dfrac{x}{6}$:
\[x^2 + 2x - 120 = 0 \iff (x-10)(x+12) = 0 \iff x = 10\text{ (vì } x > 0).\]
Chiều rộng ban đầu là $10\text{ m}$, chiều dài ban đầu là $\dfrac{360}{10} = 36\text{ m}$.
Chu vi ban đầu của mảnh vườn là:
\[C = 2(36 + 10) = \boxed{92\text{ m}}.\]
\end{loi_giai_cau}

\begin{de_cau}{Câu 6}{1,0 điểm}
Cho ba số thực dương $x, y, z$ thỏa mãn $xy + yz + zx = 1$. Tìm giá trị nhỏ nhất của biểu thức:
\[P = \dfrac{x}{\sqrt{x^2+1}} + \dfrac{y}{\sqrt{y^2+1}} + \dfrac{z}{\sqrt{z^2+1}}.\]
\end{de_cau}

\begin{loi_giai_cau}
Thay $1 = xy+yz+zx$ vào mẫu số:
\[x^2 + 1 = x^2 + xy + yz + zx = (x+y)(x+z).\]
Tương tự: $y^2 + 1 = (y+z)(y+x)$ và $z^2 + 1 = (z+x)(z+y)$.
Do đó:
\[\frac{x}{\sqrt{x^2+1}} = \frac{x}{\sqrt{(x+y)(x+z)}} = \sqrt{\frac{x}{x+y} \cdot \frac{x}{x+z}} \le \frac{1}{2}\left(\frac{x}{x+y} + \frac{x}{x+z}\right).\]
Thiết lập tương tự cho $y$ và $z$:
\[\frac{y}{\sqrt{y^2+1}} \le \frac{1}{2}\left(\frac{y}{y+z} + \frac{y}{y+x}\right),\]
\[\frac{z}{\sqrt{z^2+1}} \le \frac{1}{2}\left(\frac{z}{z+x} + \frac{z}{z+y}\right).\]
Cộng ba bất đẳng thức vế theo vế:
\[P \le \frac{1}{2}\left(\frac{x+y}{x+y} + \frac{y+z}{y+z} + \frac{z+x}{z+x}\right) = \frac{1}{2}(1 + 1 + 1) = \frac{3}{2}.\]
Dấu ``='' xảy ra khi $x = y = z = \dfrac{1}{\sqrt{3}}$.
Vậy giá trị lớn nhất là $\dfrac{3}{2}$.
\end{loi_giai_cau}
